\documentclass{jsarticle}
\usepackage[dvipdfmx,final]{graphicx}
\usepackage{amsmath}
\usepackage{amssymb}
\usepackage{chemfig}
\begin{document}
\title{The real existed theory analized with reference of books.}
\author{Masaaki Yamaguchi}
\date{}
\maketitle

      The mind of one person have with a person of mind in zone field,
   and this field also belong to be existed with the person, is persuad to
   be existed from the person with a read person.
   This permissed of mind that think of person absolutly be existed to be a person,
   and this penomone of read and a read of person, and this concerned of being thought of
   reverse of being involved with relativity of being existed with the concerned relation ship
   is a one represented person.

      読む人がいると、読まれる人がいる。この読まれる人の領域に、瞑想と思考が存在するために、
   人と人の関係性が表出して、人間の関係性という、他者との対峙ができる。
   この対峙により、読む人と読まれる人という関係から、世界の具現化ができている。
   この具現化から、誤差による人との性質の違いから、他者が形つくられている。
   この現象から、世界の契りが表現からの人の関係性から、巻き込まれることにより、
   アカシックレコードが出来ている。この思考が瞑想として、見えなくなり、この領域に潜入するために、
   思考を止めることが必要としている。その逆により、世界ができている。
   これを説いているのが、存在と無であり、機械妄想や関係妄想、思考伝播、考察伝播のある人が、
   治るためと、学ぶために、読者として、読む。それを教えてくれたのが、この本の解説本である。
   タバコの事故で家が燃えるというのは、どうでもいいです。
   Albert Eienstein博士もどうでもいいです。先生気にしないでください。その方が、すごく気にしてくれていて
   うれしすぎます。体の症状のこともわかっています。
   お姉さんは、LinkinParkで助かっていて、知っていられるから、サルトル先生を象さんと言っているのが、
   分かりました。
   この本は、アカシックレコードの一部である魔法の使い方も書いている。引き寄せの法則であり、
   瞑想が、この領域が見える、感じることの具現化の方法であるのも書いている。
   儀式にもなり、苫米地英人博士のスピード脳に生まれ変わる10分の訓練本が、
   この本の書き換えになっている。
   存在による実存が、この世には、スーパーノーマルもスーパーネイチャーもないと、存在は、対峙のみにより、
   宗教は、自己欺瞞からの表出と言われているが、読者のスコトーマを外すと、存在を実現させているのは、
   他者をイメージしている自己であり、アカシックレコードとは、歴史と未来からの両方からの蓄積であり、
   存在と無は、魔法とは、自己欺瞞からの表出と説いているが、脳科学からは、松果体の作用と、前頭葉の側頭部位、
   視覚野と、瞑想法の無行であり、現在の違う目線から存在と無を読むと、今までの目線とは、違う見方ができる本であり、
   嘔吐は、蓮覚寺先生が書き換えをしてくれていた。在籍でだけでも、いいらしい。
   手相に、仏眼相があり、旅行線があり、神様の放浪線があり、三奇紋がり、帝王線の健康線があり、
   陰徳線と陽徳線があり、向上線があり、多重知能線があり、カウンセラーの線があり、
   この世に、常能力の超能力線があるのを、証明して、現存している。
   存在と無は、今からは、違う目線で、速読法で味わえるらしい。


   RubyをVisual BasicとBasicと同じと思っている人は、芸術の賁臨になっていないです。林晴比古先生と三田玄典先生と
   柴田望洋先生を学ぶべきです。$\text{C$\sharp$とJava、C++}$も学ぶべきです。Cでの拡張ライブラリーでRubyの構造もまなぶべきです。
   ポインタをリストでファイル構造体と構造体につなげるのも学ぶべきです。それが、オブジェクト指向
   を現存させています。Rubyはかっこよすぎて、きれいです。
   私は、書き出すのが、RubyがBasicの領域からくっついて、JavaやC++の良さが
   現れていなくて、私は、Rubyの書き方が、悪すぎて、
   Basicになっていて、擬似コードがRubyになっている。間違いが無いのが、Basicになっていて、擬似コードが、
   PythonやRubyになっている。
   間違いがないスポーツができると、綺麗なのが、間違いがあると、スポーツが出来ないのが、私の能力である。
   Visual Basicもオブジェクト指向であり、モジュールをクラスにして、プロシージャがモジュールクラスになっている。
   サブクラスとは違う。Basicとは全然違う。でも、ソースコードの姿は、Pythonが一番きれいです。

\end{document}
